% Part: lambda-calculus
% Chapter: introduction
% Section: syntax

\documentclass[../../../include/open-logic-section]{subfiles}

\begin{document}

\olfileid{lam}{int}{syn}
\olsection{نحو حساب لامبدا}

نأخذ متتالية متغيرات~$x$ و$y$ و$z$،~\dots، ورموز ثوابت
$a$ و$b$ و$c$،~\dots، ثم نعرّف الحدود بالاستقراء وفق البنود:
\begin{enumerate}
\item المتغير، أيًا كان، حد.
\item والثابت، أيًا كان، حد.
\item متى كان~$M$ و$N$ حدين كان~$(MN)$ حدًا.
\item ومتى كان~$M$ حدًا و$x$ متغيرًا كان~$(\lambd[x][M])$ حدًا.
\end{enumerate}
ما جاء على الصورة~$(MN)$ نسميه \emph{تطبيقًا}، وما جاء على الصورة
$(\lambd[x][M])$ نسميه \emph{تجريدًا}.

وإذا خلا النسق من الثوابت كلها كان حساب لامبدا \emph{الخالص}.
ومعظم نظرنا إنما هو في حساب~$\lambd$ الخالص؛ فالحروف الصغيرة
كلها رموز للمتغيرات. وأما الحروف الكبيرة، مثل~$M$ و$N$،
فنرمز بها إلى حدود حساب~$\lambd$.

ونصطلح في الترميز على الأمور الآتية:

\begin{conv}
\begin{enumerate}
\item إذا أغفلنا الأقواس، رتبنا التطبيق من اليسار إلى اليمين.
  فإذا كانت~$M$ و$N$ و$P$ و$Q$ حدودًا، دل~$MNPQ$ اختصارًا
  على~$(((MN)P)Q)$.
\item ومع إغفال الأقواس نجعل لتجريد لامبدا أوسع مجال ممكن.
  فـ$\lambd[x][MNP]$ تُقرأ~$(\lambd[x][((MN)P)])$.
\item ويجوز جمع متغيرات عدة بعد لامبدا للدلالة على تتابع تجريدها؛
  فـ$\lambd[xyz][M]$ اختصار
  لـ$\lambd[x][\lambd[y][\lambd[z][M]]]$.
\end{enumerate}
\end{conv}

ومن أمثلة ذلك أن
\[
\lambd[xy][xxyx \lambd[z][xz]]
\]
اختصار للتعبير
\[
\lambd[x][\lambd[y][((((xx)y)x)(\lambd[z][(xz)]))]].
\]
فلتضبط هذه الاصطلاحات وتُحفظ. قد تشق على القارئ أول الأمر،
ثم يألفها؛ وبعد الإلف يجدها أهون من تتبع كثرة الأقواس وتشابكها.

إذا لم يختلف حدان إلا في أسماء المتغيرات المربوطة، قلنا إنهما
متكافئان بحسب~$\alpha$؛ كـ$\lambd[x][x]$ و$\lambd[y][y]$.
ويحسن أن نعدهما حدًا واحدًا؛ فقولنا إن~$M$ و$N$ واحد
لا يمنع اختلافهما بإعادة تسمية المتغيرات المربوطة.
وواقعة المتغير~$x$ مربوطة متى وقعت في مجال رابط~$\lambd[x]$،
ويربطها أقرب رابط~$\lambd[x]$ إليها؛ فإن لم تدخل في مجال أي رابط
$\lambd[x]$ فهي حرة. المثال السابق خال من المتغيرات الحرة،
وأما في
\[
(\lambd[z][yz])x
\]
فـ$y$ و$x$ حران، و$z$ مربوط.
% تصحيح مصدر (0344-I1): قُيد الربط بتطابق اسم المتغير وبالرابط الأقرب.
\end{document}
